% !TeX root = ../机器学习笔记.tex
% 由主文档通过 \input 载入；请编译主文档。
\setcounter{chapter}{16} % 保留原书章号。
\chapter{前馈神经网络与反向传播}
\label{chap:ml-networks}

神经网络把多层可训练变换复合起来，同时学习表示和预测规则。理解它的起点是函数与链式法则，而不是把训练过程视为只能由软件完成的操作。

\section{神经元、层与非线性表达}

前馈神经网络（feedforward neural network）令 $h^{(0)}=x$，对第 $l$ 层计算
\begin{equation}
 z^{(l)}=W^{(l)}h^{(l-1)}+b^{(l)},\qquad
 h^{(l)}=\phi_l(z^{(l)}).
 \label{eq:ml-network-forward}
\end{equation}
若前后层宽为 $d_{l-1},d_l$，则 $W^{(l)}\in\mathbb R^{d_l\times d_{l-1}}$、$b^{(l)}\in\mathbb R^{d_l}$，激活逐元素作用。没有非线性时，多层仿射变换仍是单个仿射变换；深度的形式并不会自动增加这类函数的表达能力。

Sigmoid 将数值压到 $(0,1)$，双曲正切压到 $(-1,1)$，修正线性单元（rectified linear unit，ReLU）为 $\max(0,z)$。非线性允许网络形成分段边界。例如二元输入的异或可写为 $\operatorname{ReLU}(x_1+x_2)-2\operatorname{ReLU}(x_1+x_2-1)$，在四个输入上恰为异或标签。存在表达式不保证训练必然找到它。

\section{前向计算与输出损失}

回归常使用线性输出和平方损失，二分类使用一个分数配 Bernoulli 交叉熵，多分类使用类别分数配 Softmax 交叉熵。输入到输出的逐层计算称前向传播。批量实现需保持样本维度与特征维度的明确约定；损失求和与取平均会改变梯度尺度，学习率必须与之匹配。

单个多分类样本若最后分数为 $z^{(L)}$、概率为 $p$、标签独热向量为 $e_y$，则损失关于分数的梯度为 $p-e_y$。计算时直接采用稳定的对数概率表达，避免先计算极小概率再取对数。模型输出分数、概率和最终类别是不同对象，不能把同一归一化重复应用。

\section{计算图与反向传播}

反向传播（backpropagation）沿计算图累积标量损失的导数。记 $\delta^{(l)}=\partial\ell/\partial z^{(l)}$，隐藏层递推为
\begin{equation}
 \delta^{(l)}=((W^{(l+1)})^\top\delta^{(l+1)})\odot\phi_l'(z^{(l)}),\qquad
 \frac{\partial\ell}{\partial W^{(l)}}=\delta^{(l)}(h^{(l-1)})^\top,\quad
 \frac{\partial\ell}{\partial b^{(l)}}=\delta^{(l)}.
 \label{eq:ml-backprop}
\end{equation}
这些式子逐项来自全微分和链式法则：$dz^{(l)}=dW^{(l)}h^{(l-1)}+db^{(l)}+W^{(l)}dh^{(l-1)}$。一个中间变量通向多条路径时，各路径贡献相加，不能只取其中一条。

自动微分（automatic differentiation）追踪实际执行的运算并应用局部导数；它不同于对表达式进行代数变形的符号微分，也不同于有限差分。反向传播只计算梯度，如何用梯度更新参数由优化器决定。

\section{初始化、梯度传播与训练稳定性}

多层反向传播涉及矩阵与激活导数连乘，反复缩小或放大就产生梯度消失或爆炸。在输入分量近似独立、中心化的粗略分析下，令权重方差与输入维数成反比有助于控制线性输出方差；对 ReLU 的尺度修正还依赖输入对称性等近似假设。这些初始化解释不是深层网络稳定训练的普遍定理。

动量更新累积梯度方向；Adam 则维护一阶矩 $m_t=\beta_1m_{t-1}+(1-\beta_1)g_t$ 和逐元素二阶矩 $v_t=\beta_2v_{t-1}+(1-\beta_2)g_t^2$，从零初始化后用
$\widehat m_t=m_t/(1-\beta_1^t)$、$\widehat v_t=v_t/(1-\beta_2^t)$
更新 $\theta_{t+1}=\theta_t-\eta_t\widehat m_t/(\sqrt{\widehat v_t}+\epsilon)$，其中 $0\leq\beta_1,\beta_2<1$、$\epsilon>0$。梯度裁剪限制一次更新的极端幅度，但不能修正错误的模型或梯度。ReLU 在零处不可微，实现通常选定一个次梯度约定。

\section{正则化与训练诊断}

随机失活（dropout）在训练时以保留概率 $q$ 独立取掩码 $m_j$，输出 $\widetilde h_j=m_jh_j/q$，条件期望等于 $h_j$；推断时使用未遮蔽激活。这一期望关系不意味着整个非线性网络的输出严格等于所有遮蔽网络的平均。批量归一化等方法还可能使训练和推断使用不同统计量，需要检查状态切换。

普通 SGD 中加入二次惩罚使参数更新含 $(1-\eta\lambda)\theta$，与相应权重衰减形式一致；自适应逐坐标缩放后，两者通常不再等价。训练循环应依次取批量、前向、计算损失、反向、更新，按独立验证选择早停时刻。调试时先尝试拟合少量样本，再检查梯度和输入尺度；连小样本都拟合不了，未必是正则化不足。

\section{例题、练习与学习检查}

\begin{example}[一个两层网络的梯度]
$f(x)=w_2\operatorname{ReLU}(w_1x)$，$x=1,y=1$，初始 $w_1=1,w_2=2$，损失为 $(f-y)^2/2$。
\end{example}
\begin{solve}
预测为二，损失为 $1/2$。在当前激活为正的区域，$\partial\ell/\partial w_2=1$、$\partial\ell/\partial w_1=2$。同时用步长 $0.1$ 更新得 $w_1=0.8,w_2=1.9$，新预测为 $1.52$，损失为 $0.1352$。两个梯度都必须由更新前的同一组参数计算。
\end{solve}

习题：验证异或表达式；推导矩阵权重梯度的维度；用有限差分核查例题并避免跨越 ReLU 折点；解释训练损失下降而验证损失上升时，延长训练与早停可能产生的不同结果。





